\chapter{JAX and XLA for PyTorch Engineers}\label{app:jax}

Every system in this book is written in PyTorch, and for good reasons: it is the language of
the open-model ecosystem, of the serving engines of \cref{ch:serving}, and of the custom
kernels of \cref{ch:cuda}. But a large share of the frontier training runs we have reasoned
about, and almost everything that runs on TPUs (\cref{app:accel}), is written in JAX
(Bradbury et al., 2018) and compiled by XLA. The two frameworks arrive at the same
collectives, the same rooflines and the same sharded-matmul cases of \cref{ch:shardcalc}
from opposite directions. PyTorch gives you an imperative program and lets you add
parallelism by hand or through DTensor; JAX gives you a pure function and a compiler that
partitions it for you, with an escape hatch when you want to write the per-device program
yourself. Reading JAX fluently is therefore not a change of allegiance but a second view of
the same machine, and the view in which the compiler's decisions are the most visible.

This appendix is written for a reader who knows PyTorch well and JAX not at all. We build the
mental model first: pure functions, parameter pytrees, explicit random keys, and tracing. We
then rewrite the tiny GPT of \cref{ch:gpt} as JAX functions and show that it reproduces the
PyTorch model's logits on the same weights to $6\cdot10^{-7}$. The second half is about
sharding: the device mesh, named shardings, the three levels of control (let \code{jit}
decide, pin an intermediate, or write the SPMD program in \code{shard\_map}), the four
matmul cases and the collective matmul of \cref{ch:shardcalc} written in JAX, and finally
how to read the StableHLO and optimised HLO that XLA produces, so that you can see which
collective it chose and why. We close with a profile and with an honest comparison of when
each framework is the better tool.

Nothing here needs an accelerator. XLA's CPU backend can be asked to pretend it has eight
devices (\code{XLA\_FLAGS=--xla\_force\_host\_platform\_device\_count=8}), and every sharded
program in this appendix and in Lab~C runs on that simulated mesh, producing the same HLO,
the same collectives and the same shardings a real eight-chip host would. The reference code
is in \ident{llmsys/jaxlab/}.

\begin{figure}[t]
\centering
\begin{tikzpicture}[
  stage/.style={draw=sysblue,fill=sysbluebg,rounded corners=2pt,minimum width=2.4cm,
    minimum height=1.15cm,align=center,font=\small},
  num/.style={circle,fill=sysblue,text=white,font=\bfseries\small,inner sep=1.5pt},
  arr/.style={-Latex,thick,sysblue!70}]
\node[stage] (s1) {Functional\\model};
\node[stage,right=0.35cm of s1] (s2) {jit, grad,\\vmap, scan};
\node[stage,right=0.35cm of s2] (s3) {XLA and\\HLO};
\node[stage,right=0.35cm of s3] (s4) {Mesh and\\named shardings};
\node[stage,below=0.8cm of s1] (s5) {Automatic\\sharding};
\node[stage,right=0.35cm of s5] (s6) {shard\_map and\\collectives};
\node[stage,right=0.35cm of s6] (s7) {Collective\\matmul};
\node[stage,right=0.35cm of s7] (s8) {Reading HLO\\and profiles};
\foreach \i in {1,...,8}{\node[num] at (s\i.north west) {\i};}
\draw[arr] (s1)--(s2); \draw[arr] (s2)--(s3); \draw[arr] (s3)--(s4);
\draw[arr] (s4.south) |- ++(0,-0.4) -| (s5.north);
\draw[arr] (s5)--(s6); \draw[arr] (s6)--(s7); \draw[arr] (s7)--(s8);
\end{tikzpicture}
\caption{Roadmap of this appendix. Stages 1--3 build the programming model and the compiler
picture; stages 4--8 redo the sharded computation of \cref{ch:shardcalc} in JAX and read
the compiler's output.}
\label{fig:jax:roadmap}
\end{figure}

% =====================================================================
\section{Why a PyTorch engineer should read JAX}\label{sec:jax:why}

Three things are easier to see in JAX than in PyTorch, and they are exactly the things this
book cares about. First, the \emph{program} is a first-class object: \code{jax.jit} traces a
Python function into an intermediate representation that you can print, lower to StableHLO,
compile, and read back as optimised HLO. In PyTorch the equivalent artefacts (an \code{fx}
graph, Inductor's generated Triton) exist but are a layer removed from the code you wrote. Second,
\emph{sharding is a property of arrays}, not of modules: an array carries a
\code{NamedSharding} that says which mesh axis each of its dimensions is split over, and the
compiler's SPMD partitioner derives every collective from those annotations. This is
\cref{ch:shardcalc}'s calculus implemented as a compiler pass, and we can check the two
against each other. Third, when the compiler's choice is not the one we want, \code{shard\_map}
lets us write the per-device program with explicit \code{psum}, \code{all\_gather} and
\code{ppermute}, while keeping \code{jit}, \code{grad} and \code{vmap} working across it.
Writing the four matmul cases once at each level is the fastest way we know to understand
what GSPMD (Xu et al., 2021) does on your behalf.

\begin{keyidea}[The same machine, two directions]
PyTorch starts from an imperative program and adds parallelism (DDP wrappers, FSDP, DTensor
placements, hand-written collectives). JAX starts from a pure function and a declaration of
where its inputs and outputs live, and a compiler fills in the collectives. Both end in the
same HLO-level program: matmuls on local shards joined by all-gathers, reduce-scatters and
all-reduces whose bytes are fixed by the shardings alone. Learning to read that program is
the skill; the framework is the notation.
\end{keyidea}

% =====================================================================
\section{The functional model}\label{sec:jax:functional}

\subsection{Pure functions and parameter pytrees}
An \code{nn.Module} bundles parameters, buffers and a \code{forward} method; calling it
mutates nothing visible but the module \emph{owns} state, and that state is what
\code{load\_state\_dict}, \code{.to(device)} and the optimizer operate on. A JAX model is a
pure function of \code{(params, inputs)} where \code{params} is a \emph{pytree}: any nesting
of Python dicts, lists and tuples whose leaves are arrays. The function has no attributes
and no device; the parameters are just a value that the caller passes in and gets new
versions of back. \Cref{lst:jax:block} is the transformer block of \cref{ch:gpt} in this
style; compare it with \code{TransformerBlock} in \ident{llmsys/gpt.py}.

\begin{lstlisting}[caption={A pre-LN transformer block as pure functions on a parameter pytree (\ident{llmsys/jaxlab/basics.py}). Linear weights are stored $(d_\text{in}, d_\text{out})$, the transpose of \code{nn.Linear.weight}.},label={lst:jax:block}]
def linear(p, x):                      # p = {"w": (in, out), "b": (out,)} or no "b"
    y = x @ p["w"]
    return y + p["b"] if "b" in p else y

def layer_norm(p, x, eps=1e-5):
    mean = x.mean(-1, keepdims=True)
    var = ((x - mean) ** 2).mean(-1, keepdims=True)
    return p["scale"] * (x - mean) / jnp.sqrt(var + eps) + p["shift"]

def causal_attention(p, x, n_heads):   # n_heads is a Python int, not a traced value
    B, T, D = x.shape
    hd = D // n_heads
    split = lambda t: t.reshape(B, T, n_heads, hd).transpose(0, 2, 1, 3)
    q, k, v = split(linear(p["q"], x)), split(linear(p["k"], x)), split(linear(p["v"], x))
    scores = q @ k.swapaxes(-1, -2) / math.sqrt(hd)
    scores = jnp.where(jnp.tril(jnp.ones((T, T), dtype=bool)), scores, -jnp.inf)
    weights = jax.nn.softmax(scores.astype(jnp.float32), axis=-1).astype(x.dtype)
    ctx = (weights @ v).transpose(0, 2, 1, 3).reshape(B, T, D)
    return linear(p["out"], ctx)

def transformer_block(p, x, n_heads):
    x = x + causal_attention(p["att"], layer_norm(p["norm1"], x), n_heads)
    return x + feed_forward(p["ff"], layer_norm(p["norm2"], x))

def gpt_forward(params, idx, n_heads):
    x = params["tok_emb"][idx] + params["pos_emb"][: idx.shape[-1]]
    for blk in params["blocks"]:       # a Python loop: unrolled when traced
        x = transformer_block(blk, x, n_heads)
    return layer_norm(params["final_norm"], x) @ params["out_head"]
\end{lstlisting}

The parameter pytree for this model is a dict with keys \code{tok\_emb}, \code{pos\_emb},
\code{blocks} (a list, one dict per layer), \code{final\_norm} and \code{out\_head}
(\cref{fig:jax:pytree}). Pytrees are what every JAX transformation operates on:
\code{jax.grad} returns a pytree of gradients with the same structure as the parameters,
\code{jax.tree.map(lambda p, g: p - lr * g, params, grads)} is an SGD step, and
\code{jax.tree.leaves(params)} flattens the whole model for counting or checkpointing. The
same structure is what a PyTorch \code{state\_dict} is, minus the module that owns it, and
\cref{sec:jax:bridge} converts one into the other mechanically.

\begin{figure}[t]
\centering
\begin{tikzpicture}[font=\scriptsize,
  box/.style={draw,rounded corners=2pt,fill=#1,align=center,minimum height=0.55cm},
  arr/.style={-Latex,thick}]
% PyTorch side
\node[font=\small\bfseries] at (0,3.2) {PyTorch};
\node[box=sysblue!15,minimum width=4.6cm,minimum height=2.6cm] (mod) at (0,1.4) {};
\node[anchor=north west] at (mod.north west) {\texttt{GPTModel(nn.Module)}};
\node[box=keygreen!25,minimum width=3.9cm] at (0,1.9) {\texttt{.tok\_emb .pos\_emb .trf\_blocks[i]}};
\node[box=keygreen!25,minimum width=3.9cm] at (0,1.25) {\texttt{.final\_norm .out\_head} (owned state)};
\node[box=napkin!25,minimum width=3.9cm] at (0,0.55) {\texttt{forward(idx)} reads \texttt{self.*}};
\node[below=0.1cm of mod,align=center,text width=4.6cm] {state lives inside the object; \texttt{.to()}, optimizer and DDP act on it};
% JAX side
\node[font=\small\bfseries] at (7.2,3.2) {JAX};
\node[box=keygreen!25,minimum width=3.3cm,minimum height=2.6cm] (tree) at (5.6,1.4) {};
\node[anchor=north west] at (tree.north west) {\texttt{params} (pytree)};
\node[box=white,minimum width=2.7cm] at (5.6,1.95) {\texttt{tok\_emb pos\_emb}};
\node[box=white,minimum width=2.7cm] at (5.6,1.3) {\texttt{blocks: [\{norm1,att,...\}]}};
\node[box=white,minimum width=2.7cm] at (5.6,0.65) {\texttt{final\_norm out\_head}};
\node[box=napkin!25,minimum width=3.0cm,minimum height=1.2cm,align=center] (fn) at (9.3,1.4) {\texttt{gpt\_forward(}\\\texttt{params, idx, H)}\\pure function};
\draw[arr] (tree.east) -- (fn.west) node[midway,above] {argument};
\node[below=0.1cm of tree.south,align=center,text width=3.3cm,xshift=1.8cm] {state is a value; \texttt{grad}, \texttt{jit}, \texttt{vmap} and sharding act on the function and the pytree};
\draw[arr,labpurple] (mod.east) -- (tree.west) node[midway,above,align=center] {\texttt{state\_dict}\\(transpose $W$)};
\end{tikzpicture}
\caption{A module owns its parameters; a JAX model receives them. The pytree has the same
shape as a \code{state\_dict}, which is why converting between the two is mechanical.}
\label{fig:jax:pytree}
\end{figure}

\subsection{Immutability}
JAX arrays cannot be modified in place. \code{x[i] = v} is an error; \code{x.at[i].set(v)}
returns a new array with the update applied. Under \code{jit} this costs nothing, because XLA
owns the buffers and turns the functional update into an in-place one when the old value is
dead (\emph{buffer donation} lets you promise this across a \code{jit} boundary too, with
\code{donate\_argnums}). What you lose is the PyTorch idiom of preallocating a KV cache and
writing into it from the outside; what you gain is that every function is safe to retrace,
reorder, differentiate and partition, because no operation has a hidden effect on anything
else. The compiler's freedom to fuse, schedule and shard rests entirely on this guarantee.

\subsection{Explicit randomness}
PyTorch has a global generator: \code{torch.manual\_seed(0)} and every subsequent
\code{torch.randn} advances hidden state. A pure function cannot do that, so JAX threads
randomness through explicitly. A \emph{key} is a value, \code{jax.random.key(0)};
\code{jax.random.split(key, n)} derives $n$ independent keys; each sampling call consumes a
key and the same key always gives the same numbers (\cref{lst:jax:rng}). The discipline looks
pedantic until you shard a model: with explicit keys, dropout on device 3 uses a key that is
a pure function of (step, layer, device), so a run is reproducible across any number of
devices and any re-partitioning, which is not true of a per-process global generator.

\begin{lstlisting}[caption={Explicit PRNG keys: dropout as a pure function of its key.},label={lst:jax:rng}]
def dropout(key, x, rate):
    keep = jax.random.bernoulli(key, 1.0 - rate, x.shape)
    return jnp.where(keep, x / (1.0 - rate), 0.0)

key = jax.random.key(0)
k_init, k_drop = jax.random.split(key)          # never reuse a key for two purposes
params = init_gpt_params(k_init, cfg)
h = dropout(k_drop, h, 0.1)                     # same k_drop -> same mask, on any device
\end{lstlisting}

\subsection{A translation table}
\Cref{tab:jax:mental} collects the correspondences we use in the rest of the appendix. The
right column is deliberately short: most of JAX's API is NumPy's (\code{jax.numpy} is a
drop-in for the subset you need), and the framework-specific part is a handful of
transformations and the sharding types.

\begin{table}[h]\centering\small
\begin{tabularx}{\linewidth}{>{\raggedright\arraybackslash}X>{\raggedright\arraybackslash}X}
\toprule
PyTorch & JAX \\
\midrule
\code{nn.Module} holding parameters & pure function + parameter pytree passed in \\
\code{torch.manual\_seed}, global RNG & explicit \code{jax.random.key}, split per use \\
in-place ops (\code{x.add\_()}, \code{x[i] = v}) & immutable arrays: \code{x.at[i].set(v)} returns a new array \\
\code{loss.backward(); p.grad} & \code{grads = jax.grad(loss\_fn)(params, ...)} \\
\code{torch.compile} (Dynamo bytecode capture, guards) & \code{jax.jit} (tracing to jaxpr, one trace per shape signature) \\
\code{torch.func.vmap} & \code{jax.vmap} \\
\code{for layer in layers: x = layer(x)} & \code{jax.lax.scan} over stacked layer parameters \\
\code{DeviceMesh} + DTensor \code{Shard(i)}/\code{Replicate()} & \code{Mesh} + \code{NamedSharding(mesh, PartitionSpec(...))} \\
\code{dist.all\_reduce} written by you & \code{shard\_map} + \code{jax.lax.psum}, or let \code{jit} insert it \\
\code{torch.cuda.synchronize()} & \code{jax.block\_until\_ready(x)} (dispatch is asynchronous) \\
\code{torch.fx} graph / Inductor Triton & jaxpr / StableHLO / optimised HLO (\cref{sec:jax:xla}) \\
\bottomrule
\end{tabularx}
\caption{The PyTorch-to-JAX mental model used in this appendix (\code{basics.MENTAL\_MODEL}
prints it).}
\label{tab:jax:mental}
\end{table}

\begin{systemslens}[Purity is what the compiler is paid in]
Every optimisation XLA performs---fusing a chain of elementwise ops into one kernel, choosing
memory layouts, deciding which device computes which block of a matmul, overlapping a
collective with the matmul that follows it---requires knowing that reordering two operations
cannot change the result. In an imperative framework the compiler has to \emph{prove} this
(Dynamo guards, aliasing analysis, graph breaks when it cannot); in JAX the programming model
guarantees it. That is the trade: you give up in-place updates and global state, and in
return the whole program, including its parallelisation, becomes the compiler's to
rearrange.
\end{systemslens}

% =====================================================================
\section{Transformations: jit, grad, vmap, scan}\label{sec:jax:transforms}

\subsection{jit: tracing, not capturing}
\code{jax.jit(f)} returns a function that, on its first call with arrays of a given shape and
dtype, runs \code{f} once with \emph{tracers} in place of the arrays. Every \code{jnp}
operation on a tracer records itself into a \emph{jaxpr}, a small typed IR
(\cref{lst:jax:jaxpr}). Python control flow executes during tracing: an \code{if} on a
Python value picks a branch at trace time, a \code{for} loop over a Python range is unrolled,
and a Python \code{int} like \code{n\_heads} is a compile-time constant. Anything that depends
on the \emph{value} of a traced array (\code{if x.sum() > 0}) is an error, because the
tracer has a shape but no value; \code{jax.lax.cond} and \code{jax.lax.while\_loop} are the
traced alternatives. The jaxpr is then lowered to StableHLO, compiled by XLA, and cached
under the \emph{signature}: the pytree structure, shapes and dtypes of the arguments plus any
\code{static\_argnums}. A later call with the same signature skips Python entirely and
dispatches the compiled executable.

\begin{lstlisting}[style=plain,caption={The jaxpr of \code{jax.nn.gelu(x @ w)} for $x\in\R^{8\times16}$, $w\in\R^{16\times8}$ (\code{jax.make\_jaxpr}). Each line is one primitive with typed inputs and outputs.},label={lst:jax:jaxpr}]
{ lambda ; a:f32[8,16] b:f32[16,8]. let
    c:f32[8,8] = dot_general[dimension_numbers=(([1], [0]), ([], []))] a b
    d:f32[8,8] = integer_pow[y=3] c
    e:f32[8,8] = mul 0.044714998453855515:f32[] d
    f:f32[8,8] = add c e
    g:f32[8,8] = mul 0.7978845834732056:f32[] f
    h:f32[8,8] = tanh g
    i:f32[8,8] = add 1.0:f32[] h
    j:f32[8,8] = mul 0.5:f32[] i
    k:f32[8,8] = mul c j
  in (k,) }
\end{lstlisting}

This is a different mechanism from \code{torch.compile}. Dynamo captures a graph by
symbolically executing Python \emph{bytecode}, installing guards on the values it specialised
on, and falling back to eager Python at a graph break; it can handle data-dependent control
flow by splitting the graph around it, and with dynamic shapes it compiles a symbolic program
once for many sizes. JAX's tracing is simpler and stricter: there are no graph breaks (a
function either traces or raises), and shapes are always static, so a new shape means a new
trace and a new compile (\cref{fig:jax:tracing}). \Cref{lst:jax:traces} makes the cache
visible: the Python body runs once per new shape and never again for a shape it has seen.

\begin{figure}[t]
\centering
\begin{tikzpicture}[font=\scriptsize,
  box/.style={draw,rounded corners=2pt,fill=#1,align=center,minimum height=0.6cm,minimum width=1.9cm},
  arr/.style={-Latex,thick}]
\node[box=sysblue!15] (call) at (0,0) {\texttt{f(x)} call};
\node[box=napkin!25] (sig) at (2.6,0) {signature:\\shapes, dtypes,\\static args};
\node[draw,diamond,aspect=2,fill=white,inner sep=1pt] (hit) at (5.4,0) {cached?};
\node[box=keygreen!25] (run) at (8.4,0.9) {run executable\\(no Python)};
\node[box=warnred!15] (trace) at (8.4,-0.9) {trace body $\to$ jaxpr\\$\to$ StableHLO $\to$ XLA};
\node[box=fmemory!30] (cache) at (11.4,-0.9) {cache[sig]\\= executable};
\draw[arr] (call)--(sig); \draw[arr] (sig)--(hit);
\draw[arr] (hit) |- node[pos=0.3,left] {yes} (run);
\draw[arr] (hit) |- node[pos=0.3,left] {no} (trace);
\draw[arr] (trace)--(cache); \draw[arr] (cache) |- ++(0,1.8) -| (run);
\node[align=left,anchor=north west,text width=11.5cm] at (-1.0,-1.9)
 {Dynamo instead: symbolic execution of bytecode $\to$ FX graph with \emph{guards}; a failed guard recompiles,
  an unsupported construct splits the graph and runs that piece eagerly. JAX has no graph breaks and no symbolic shapes:
  a shape it has not seen is a new program.};
\end{tikzpicture}
\caption{What a call to a jitted function does. The cost of the ``no'' branch is the
shape-specialisation cost of \cref{sec:jax:xla}; the ``yes'' branch is a hash lookup and an
asynchronous dispatch.}
\label{fig:jax:tracing}
\end{figure}

\begin{lstlisting}[caption={Counting traces: the Python body runs only on a cache miss (\code{basics.count\_traces}).},label={lst:jax:traces}]
def count_traces(fn, inputs):
    calls = [0]
    def counted(*args):
        calls[0] += 1                       # Python runs only while tracing
        return fn(*args)
    jitted = jax.jit(counted)
    history = []
    for args in inputs:
        jax.block_until_ready(jitted(*args))
        history.append(calls[0])
    return history

count_traces(lambda x: x * 2, [(jnp.ones(3),), (jnp.ones(3),), (jnp.ones(4),)])
# -> [1, 1, 2]: same shape reuses the executable, a new shape retraces and recompiles
\end{lstlisting}

\begin{pitfall}[Python values versus traced values]
Inside a jitted function, \code{x.shape} is a tuple of Python ints and may drive control
flow; \code{x[0]} is a tracer and may not. Passing a model hyper-parameter such as
\code{n\_heads} as an ordinary argument makes it a traced \code{int32} scalar, and
\code{D // n\_heads} then fails to produce a static shape for the reshape. Mark it
\code{static\_argnums} (or close over it), and remember that every distinct static value is
a separate compile. The same rule applies to \code{jax.lax.scan} lengths, mesh shapes and
\code{PartitionSpec}s: all are compile-time constants.
\end{pitfall}

\subsection{grad and value\_and\_grad}
\code{jax.grad(loss\_fn)} is a function that returns $\partial\text{loss}/\partial(\text{first
argument})$ as a pytree shaped like that argument. There is no \code{.backward()}, no
\code{.grad} attribute and no \code{zero\_grad()}: gradients are return values. JAX
implements this by \emph{tracing} the function to a jaxpr, applying a forward-mode rule
(JVP) to every primitive and then \emph{transposing} the resulting linear program to obtain
reverse mode (Frostig et al., 2018). The transposition step is the one to remember: it is why,
in \cref{sec:jax:shardmap}, the gradient of an \code{all\_gather} comes out as a
\code{psum\_scatter} and the gradient of a \code{ppermute} as the reverse permutation, with no
code of ours.

\begin{lstlisting}[caption={Loss and gradient for the JAX GPT. \code{value\_and\_grad} returns both in one traced program.},label={lst:jax:grad}]
def cross_entropy(logits, targets):
    logp = jax.nn.log_softmax(logits.astype(jnp.float32), axis=-1)
    return -jnp.take_along_axis(logp, targets[..., None], axis=-1)[..., 0].mean()

def loss_fn(params, idx, targets, n_heads):
    return cross_entropy(gpt_forward(params, idx, n_heads), targets)

step = jax.jit(jax.value_and_grad(loss_fn), static_argnums=3)
loss, grads = step(params, idx, targets, 4)          # grads: same pytree as params
params = jax.tree.map(lambda p, g: p - 1e-3 * g, params, grads)   # SGD, functionally
\end{lstlisting}

\subsection{vmap: batching as a transformation}
\code{jax.vmap(f, in\_axes)} turns a function of one example into a function of a batch by
adding a batch dimension to every primitive in its jaxpr; the result is a single batched
program, not a loop. Combined with \code{grad} it gives per-example gradients
(\cref{lst:jax:vmap}), the primitive behind gradient clipping per sample, differential
privacy, and the influence computations of \cref{ch:data}; in PyTorch the same thing is
\code{torch.func.vmap(torch.func.grad(...))} and is a more recent addition built on the same
idea. \code{in\_axes=(None, 0, 0)} says: do not map over the parameters, map the two token
arrays along their first axis.

\begin{lstlisting}[caption={Per-example gradients: \code{vmap} of \code{grad} of a one-example loss (\code{basics.per\_example\_grads}).},label={lst:jax:vmap}]
def per_example_grads(params, idx, targets, n_heads):
    one = lambda p, i, t: loss_fn(p, i[None], t[None], n_heads)    # (T,) -> (1, T)
    return jax.vmap(jax.grad(one), in_axes=(None, 0, 0))(params, idx, targets)
# every leaf of the result has a leading axis of size B; one batched program runs
\end{lstlisting}

\subsection{scan: a loop the compiler can see}\label{sec:jax:scan}
The Python loop in \cref{lst:jax:block} is unrolled during tracing: a 32-layer model becomes
32 copies of the block in the jaxpr, in the StableHLO and in the compiled program, and compile
time grows with depth. \code{jax.lax.scan} traces the body \emph{once} and emits a loop whose
carried state is the activation and whose per-iteration input is layer $\ell$'s slice of a
\emph{stacked} parameter pytree (\cref{lst:jax:scan}). Production JAX LLMs store their layers
stacked for exactly this reason, and it is why a \code{state\_dict} converted from PyTorch
needs a \code{stack\_blocks} step before it fits such a model.

\begin{lstlisting}[caption={The layer loop as \code{lax.scan}: one traced body, stacked parameters (\code{basics.gpt\_forward\_scan}).},label={lst:jax:scan}]
def stack_blocks(blocks):            # list of per-layer pytrees -> one pytree with a layer axis
    return jax.tree.map(lambda *leaves: jnp.stack(leaves), *blocks)

def gpt_forward_scan(params, idx, n_heads):
    x = params["tok_emb"][idx] + params["pos_emb"][: idx.shape[-1]]
    def body(x, blk):                # blk is layer l's slice of the stacked pytree
        return transformer_block(blk, x, n_heads), None
    x, _ = jax.lax.scan(body, x, params["blocks"])      # params["blocks"] is stacked
    return layer_norm(params["final_norm"], x) @ params["out_head"]
\end{lstlisting}

\begin{napkinmath}[Program size with and without scan]
For the tiny GPT (\code{GPT\_CONFIG\_TINY}, $\dmodel=64$) on the CPU backend we measured the
lowered program for 2, 8 and 16 layers. Unrolled: 317, 1067 and 2067 lines of StableHLO and
1086, 3480 and 6672 lines of optimised HLO, growing by about 125 StableHLO and 400 HLO lines
per layer. With \code{scan}: 253 and 865 lines at every depth. Compile time followed the same
pattern (0.31, 0.38, 0.59\,s unrolled versus a flat 0.73\,s with scan: the loop costs a
little fixed overhead and the unrolled program only overtakes it beyond about 20 layers at
this size). At production width the unrolled program is tens of thousands of HLO
instructions and compiles for minutes; scan keeps it at one layer's worth. The price is that
the compiler cannot specialise layer $\ell$ differently from layer $\ell'$, which matters
for the first and last layers of pipelines and for models with heterogeneous layers
(the hybrids of \cref{ch:attnalt}): those are written as a scan over the homogeneous
middle plus explicit first and last blocks.
\end{napkinmath}

\subsection{The PyTorch bridge}\label{sec:jax:bridge}
To be sure that \cref{lst:jax:block} is the \emph{same} model as \ident{llmsys/gpt.py} and
not merely a similar one, we convert the PyTorch model's \code{state\_dict} into the pytree
and compare logits (\cref{lst:jax:bridge}). The conversion is a renaming plus one transpose
per linear layer; the tensors cross over as NumPy arrays, so \code{torch} is never imported
on the JAX side. On the tiny configuration the largest logit difference is
$6\cdot10^{-7}$ in fp32, which is the accumulated rounding difference of two libraries'
matmul orderings and nothing else.

\begin{lstlisting}[caption={Converting an \code{llmsys.gpt.GPTModel} state\_dict to the pytree of \cref{lst:jax:block} (\code{basics.torch\_state\_dict\_to\_jax}, abridged).},label={lst:jax:bridge}]
def torch_state_dict_to_jax(sd):
    arr = lambda name: jnp.asarray(sd[name].detach().cpu().float().numpy())
    def lin(prefix):                              # nn.Linear.weight is (out, in)
        p = {"w": arr(prefix + ".weight").T}
        if prefix + ".bias" in sd:
            p["b"] = arr(prefix + ".bias")
        return p
    norm = lambda prefix: {"scale": arr(prefix + ".scale"), "shift": arr(prefix + ".shift")}
    blocks = [{"norm1": norm(f"trf_blocks.{i}.norm1"),
               "att": {"q": lin(f"trf_blocks.{i}.att.W_query"), ...},
               "norm2": norm(f"trf_blocks.{i}.norm2"),
               "ff": {"fc1": lin(f"trf_blocks.{i}.ff.layers.0"),
                      "fc2": lin(f"trf_blocks.{i}.ff.layers.2")}} for i in range(n_layers)]
    return {"tok_emb": arr("tok_emb.weight"), "pos_emb": arr("pos_emb.weight"),
            "blocks": blocks, "final_norm": norm("final_norm"),
            "out_head": arr("out_head.weight").T}

params = torch_state_dict_to_jax(model.state_dict())
jax_logits = gpt_forward(params, jnp.asarray(idx), model.cfg["n_heads"])
# max |jax_logits - model(idx)| = 6e-7 on GPT_CONFIG_TINY
\end{lstlisting}

% =====================================================================
\section{XLA: from jaxpr to machine code}\label{sec:jax:xla}

\subsection{The pipeline}
XLA (Accelerated Linear Algebra) is a domain-specific compiler whose input is HLO
(High-Level Operations), a small graph language of about a hundred ops: elementwise
arithmetic, \code{dot}, \code{reduce}, \code{broadcast}, \code{reshape}, \code{dynamic-slice},
\code{while}, \code{conditional}, the collectives, and a few more. JAX lowers a jaxpr to
\emph{StableHLO}, the versioned MLIR dialect that is XLA's public input format, and hands it
to XLA, which runs a long sequence of passes (\cref{fig:jax:pipeline}) and emits an
executable for the target: LLVM IR for CPU, PTX for NVIDIA GPUs, and the proprietary
backend for TPUs. Two of those passes are the subject of the rest of this appendix: the
\emph{SPMD partitioner}, which turns one program with sharding annotations into the
per-device program with collectives, and \emph{fusion}. Two more decide most of the
single-device performance: \emph{layout assignment} and \emph{scheduling with buffer
assignment}.

\begin{figure}[t]
\centering
\resizebox{\linewidth}{!}{%
\begin{tikzpicture}[font=\scriptsize,
  st/.style={draw,rounded corners=2pt,fill=#1,align=center,minimum height=0.9cm,minimum width=2.0cm},
  arr/.style={-Latex,thick}]
\node[st=sysblue!15] (py) at (0,0) {Python\\\texttt{f(params, x)}};
\node[st=sysblue!15] (jp) at (2.5,0) {jaxpr\\(trace)};
\node[st=keygreen!25] (sh) at (5.0,0) {StableHLO\\\texttt{.lower().as\_text()}};
\node[st=napkin!25,minimum width=7.6cm] (xla) at (10.6,0) {XLA passes: SPMD partitioner $\to$ algebraic simplifier $\to$\\
  fusion $\to$ layout assignment $\to$ scheduling $\to$ buffer assignment};
\node[st=keygreen!25] (hlo) at (16.2,0) {optimised HLO\\\texttt{.compile().as\_text()}};
\node[st=fmemory!30] (exe) at (18.9,0) {LLVM / PTX /\\TPU executable};
\draw[arr] (py)--(jp); \draw[arr] (jp)--(sh); \draw[arr] (sh)--(xla); \draw[arr] (xla)--(hlo); \draw[arr] (hlo)--(exe);
\node[anchor=north,align=center,text width=4.6cm] at (3.75,-0.7) {one trace per shape signature};
\node[anchor=north,align=center,text width=7.4cm] at (10.6,-0.7) {all shapes static: every buffer size and every collective's byte count is known here};
\end{tikzpicture}}
\caption{From a Python function to an executable. The two texts a JAX user can read are
StableHLO (what JAX asked for) and optimised HLO (what XLA decided to run); the lab dumps
both for every sharded matmul case.}
\label{fig:jax:pipeline}
\end{figure}

\subsection{Reading StableHLO and HLO}
\Cref{lst:jax:stablehlo} is the StableHLO of \code{jax.nn.gelu(x @ w)}: one op per jaxpr
primitive, in MLIR syntax, with scalar constants broadcast explicitly. \Cref{lst:jax:hlo} is
the optimised HLO of the same program. Three things happened. The ten elementwise ops became
\emph{one} loop fusion (\code): a single kernel that reads the matmul output once
and writes the GELU once, which is exactly the fusion we wrote by hand in Triton in
\cref{app:triton}. The \code{dot} became a \code{kCustom} fusion that calls the backend's
matmul library (on CPU, XLA's own tiled kernels; on GPU, cuBLAS or a Triton GEMM). And every
array acquired a \emph{layout}: \code{f32[8,8]\{1,0\}} means an $8\times8$ array whose
dimension 1 is minor-most (row-major). Layout assignment chooses these per instruction so
that the matmul library gets the layout it wants and inserts transposes only where it must;
on TPU this is where the $(8,128)$ tiling of \cref{app:accel} enters.

\begin{lstlisting}[style=plain,caption={StableHLO for \code{gelu(x @ w)} (\code{hlo.stablehlo\_text}, abridged). Every jaxpr primitive is one op; nothing has been fused yet.},label={lst:jax:stablehlo}]
func.func public @main(%arg0: tensor<8x16xf32>, %arg1: tensor<16x8xf32>) -> tensor<8x8xf32> {
  %0 = stablehlo.dot_general %arg0, %arg1, contracting_dims = [1] x [0]
  %1 = stablehlo.multiply %0, %0 : tensor<8x8xf32>
  %2 = stablehlo.multiply %1, %0 : tensor<8x8xf32>
  %cst = stablehlo.constant dense<4.471500e-02> : tensor<f32>
  %3 = stablehlo.broadcast_in_dim %cst, dims = [] : (tensor<f32>) -> tensor<8x8xf32>
  %4 = stablehlo.multiply %3, %2 : tensor<8x8xf32>
  %5 = stablehlo.add %0, %4 : tensor<8x8xf32>
  ...                                  # tanh, 1 + ., 0.5 * ., x * .
  return %13 : tensor<8x8xf32>
}
\end{lstlisting}

\begin{lstlisting}[style=plain,caption={Optimised HLO for the same program (\code{hlo.compiled\_text}, metadata stripped). Two fusions: the matmul and one loop fusion for the whole GELU.},label={lst:jax:hlo}]
%fused_computation.1 (param_0.1: f32[8,8]) -> f32[8,8] {
  %param_0.1 = f32[8,8]{1,0} parameter(0)
  %integer_pow.1 = f32[8,8]{1,0} multiply(%param_0.1, %param_0.1)
  %integer_pow.0 = f32[8,8]{1,0} multiply(%integer_pow.1, %param_0.1)
  ...                                  # the remaining elementwise ops
  ROOT %mul.0 = f32[8,8]{1,0} multiply(%param_0.1, %mul.1)
}

ENTRY %main.1 (x.1: f32[8,16], w.1: f32[16,8]) -> f32[8,8] {
  %x.1 = f32[8,16]{1,0} parameter(0)
  %w.1 = f32[16,8]{1,0} parameter(1)
  %ynn_fusion = f32[8,8]{1,0} fusion(%x.1, %w.1), kind=kCustom, calls=%fused_computation
  ROOT %multiply_multiply_fusion = f32[8,8]{1,0} fusion(%ynn_fusion), kind=kLoop,
        calls=%fused_computation.1
}
\end{lstlisting}

The HLO grammar is regular enough to parse with a regular expression, which is what
\code{hlo.count\_fusions} and \code{hlo.count\_collectives} do: every instruction is
\code{\%name = type op(operands), attributes}, and the \code{ENTRY} computation is the
program's top level. Fusions come in four kinds. \code{kLoop} is a fused elementwise loop;
\code{kInput} is a reduction (softmax's row max and sum, a norm's mean of squares) with its
producers fused in; \code{kOutput} fuses an elementwise epilogue into a matmul or
convolution (the bias-and-activation epilogue of \cref{app:triton}); \code{kCustom} calls a
library or a hand-written kernel. Reading the fusion list of a transformer layer tells you
in seconds whether the norm, the activation and the residual add became one kernel each or
one kernel together.

\subsection{The shape-specialisation trade-off}
XLA compiles a program for \emph{exactly} the shapes it was traced with. This is the source
of both its strength and its most-cited weakness. With every shape known, every buffer's
size is known, so XLA can assign all temporaries to a single preplanned arena (buffer
assignment), schedule instructions to minimise peak memory, decide fusion profitability
exactly, and emit collectives whose byte counts are compile-time constants. The price is one
compile per distinct shape, and compiles of large programs take minutes.

\begin{napkinmath}[How many shapes can you afford?]
Suppose a training-step program compiles in \SI{90}{\second} and a step runs in
\SI{2}{\second}. A single shape amortises its compile over the first 45 steps of a run that
has $10^5$ of them: negligible. Now consider serving with variable prompt lengths. Compiling
one program per length from 1 to 4096 is $4096\times\SI{90}{\second}\approx\SI{100}{\hour}$,
which is absurd; compiling for 13 power-of-two \emph{buckets} (1, 2, 4, \dots, 4096) and
padding each prompt up to the next bucket is \SI{20}{\minute} of compile and at most
$2\times$ wasted prefill compute, $1.4\times$ on average. This is exactly the bucketing that
TPU serving systems do, and it is why the paged KV cache of \cref{ch:serving}, whose block
table has a data-dependent shape, needs a fixed maximum number of blocks per sequence in a
JAX implementation. The compile cache can be persisted to disk across processes
(\code{jax.config.update("jax\_compilation\_cache\_dir", ...)}), which turns the \SI{20}{\minute}
into a one-off.
\end{napkinmath}

\begin{pitfall}[Dynamic shapes are a redesign, not a flag]
A PyTorch program that grows a KV cache by concatenation, drops finished sequences from a
batch, or routes a variable number of tokens to each expert will retrace and recompile on
every step under \code{jit}. The fixes are the ones \cref{ch:inference} and
\cref{ch:moe} already motivated for performance reasons: preallocate the cache at maximum
length and write into it with \code{dynamic\_update\_slice}; keep the batch at a fixed size
with masked padding slots; give every expert a fixed capacity and drop or pad. XLA simply
refuses to let you postpone these decisions.
\end{pitfall}

\begin{systemslens}[Static shapes buy a memory plan]
The reason a JAX training run reports a single peak-memory number at compile time, and a
PyTorch run has to be probed with \code{max\_memory\_allocated}, is buffer assignment: with
every shape fixed, XLA lays out all temporaries of the step in one arena and knows the
high-water mark before the first step executes. \code{compiled.memory\_analysis()} returns it
(\code{hlo.cost\_summary} prints it alongside XLA's FLOP and byte estimates). The same
property is what lets XLA schedule collectives to overlap with compute without a runtime
allocator getting in the way, which we return to in \cref{sec:jax:cmm}.
\end{systemslens}

